\section{Paid norm divisors and two-step compatibility}
\label{tt-section}
\noindent\textbf{Status.} This section contains ordinary mathematical proofs,
checked by separate research agents, and separately identified exact finite
computations. It proves sufficient criteria for an $abc$ counterexample
family, but does not prove the existence premises of those criteria.
Arbitrary depth at one or two primes is not identified with compression of
an entire norm. No full Lean verification of this section is asserted.

\subsection{From a shared representation to the complete norm radical}
Let $a,b$ be coprime positive integers, $c=a+b$, and
\[
 z=a+b\zeta,\quad \zeta^2-\zeta+1=0,\quad
 M=N(z)=a^2+ab+b^2,\quad
 t=\log c,\quad r=\log\operatorname{rad}(abc),\quad q=t/r.
\]
Suppose there is an actual oriented representation
\begin{equation}\label{tt-profile}
 z=u(1+\zeta)^{e_0}v w^g,\qquad
 e_0\in\{0,1\},\quad g\ge1,\quad V=N(v)\ge1,\quad Q=N(w)\ge7.
\end{equation}
Set $h_v=\max(1,\log V)$, $\lambda=h_v/g$,
$\rho=\lambda\log(4+g)$, and $s=\log\operatorname{rad}(M)$.
Here $\lambda$ is a positive real compression ratio, not the complex
quotient used earlier in the logarithmic-form argument. The norm identity
and height comparison give
\begin{equation}\label{tt-norm-budget}
 M=3^{e_0}VQ^g,\qquad M\le c^2,\qquad t\ge g\log7/2.
\end{equation}

\begin{theorem}[Complete norm radical bound]\label{tt-norm-radical}
Every representation \eqref{tt-profile} satisfies
\[
 \operatorname{rad}(M)\le3^{e_0}VQ,
 \qquad
 s\le\frac{\log M}{g}
   +\left(1-\frac1g\right)(\log V+e_0\log3)
 \le\frac{2t}{g}+\log V+\log3.
\]
Consequently, with
\[
 K_*=2+\frac{2(1+\log3)}{\log7},
\]
one has
\begin{equation}\label{tt-small-radical}
 0<\frac{s}{t}\le K_*\lambda
 =\frac{K_*\rho}{\log(4+g)}.
\end{equation}
If $e_0=0$, the constant can be replaced by $2+2/\log7$.
\end{theorem}
\begin{proof}
Radicals are invariant under positive powers and submultiplicative, so
\[
 \operatorname{rad}(3^{e_0}VQ^g)
 \le3^{e_0}\operatorname{rad}(V)\operatorname{rad}(Q)
 \le3^{e_0}VQ.
\]
Overlapping support of $V,Q$ only improves this bound. Write
$x=\log V+e_0\log3$. The identity $\log M=x+g\log Q$ rewrites
$s\le x+\log Q$ as the first logarithmic inequality. Then use $g\ge1$
and \eqref{tt-norm-budget}. Finally
\[
 \frac{s}{t}\le\frac2g+
 \frac{2(h_v+\log3)}{g\log7}
 \le K_*\frac{h_v}{g},
\]
because $h_v\ge1$. Omitting the ramified factor proves the last assertion.
\end{proof}

This theorem controls the norm radical, not $\operatorname{rad}(abc)$.
Neither it nor the established balance theorem implies that the input quality
exceeds one half.

\subsection{Exact quality transport and its quantified conditional criterion}
The primitive mixed transform is
\[
 \mathcal T(a,b,c)=(ab,M,c^2).
\]
The norm-boundary gcd identity gives
$\gcd(M,abc)=1$, hence its logarithmic height and radical are
$t'=2t$ and $r'=r+s$. Therefore
\begin{equation}\label{tt-quality}
 q'=\frac{2q}{1+q(s/t)},\qquad
 \frac1{q'}=\frac1{2q}+\frac{s}{2t},\qquad
 \frac1{2q}\le\frac1{q'}\le\frac1{2q}+\frac{K_*}2\lambda.
\end{equation}
These identities follow by division of $t'=2t$ by $r'=r+s$.
In particular, along a sequence with $\rho\to0$,
$1/q'-1/(2q)\to0$. Claiming $q'/q\to2$ additionally requires
$q(s/t)\to0$, for example bounded input qualities; it is not asserted
without that extra condition.

\begin{theorem}[Half-quality sufficient criterion]\label{tt-half-quality}
Suppose a sequence of actual primitive triples admits
\eqref{tt-profile}, has $\rho_n\to0$, and satisfies
$q_n\ge q_0>1/2$ eventually for one fixed $q_0$. Then, for one fixed
$\varepsilon>0$,
\[
 \frac{c_n^2}
 {\bigl(\operatorname{rad}(a_nb_nc_n)\operatorname{rad}(M_n)\bigr)^{1+\varepsilon}}
 \longrightarrow\infty.
\]
Thus such a sequence would disprove standard $abc$.
\end{theorem}
\begin{proof}
Choose once and for all $0<\varepsilon<2q_0-1$, and put
$d=2-(1+\varepsilon)/q_0>0$. The output excess is exactly
\[
 E'_\varepsilon
 =\left(2-\frac{1+\varepsilon}{q}
          -(1+\varepsilon)\frac{s}{t}\right)t.
\]
Equation~\eqref{tt-small-radical} implies $s_n/t_n\to0$, so eventually
$E'_{\varepsilon,n}\ge dt_n/2$. Also $\rho_n\ge1/g_n$, hence
$g_n\to\infty$, and \eqref{tt-norm-budget} gives $t_n\to\infty$.
The fixed-epsilon excess therefore tends to infinity. Exponentiating defeats
every uniform constant. The outputs remain primitive by the gcd identity.
\end{proof}

The same statement follows from $\limsup q_n>1/2$ by passage to a
subsequence. No such sequence is constructed here. The previously constructed
small-$\rho$ families are proved to be balanced; that does not prove this
quality premise. The elementary bound
$\operatorname{rad}(abc)\le abc\le c^3/4$ supplies only an asymptotic
one-third lower bound for quality.

\subsection{Fresh divisors over several steps}
Iterate the actual transform by
\[
 (a_{j+1},b_{j+1},c_{j+1})=(a_jb_j,M_j,c_j^2),\qquad
 M_j=a_j^2+a_jb_j+b_j^2.
\]
Every new $M_j$ is coprime to the entire accumulated support. Thus the
$M_j$ are pairwise coprime and coprime to the initial boundary, and induction
gives
\begin{equation}\label{tt-telescope}
 t_n=2^nt_0,\quad
 R_n=R_0\prod_{j<n}\operatorname{rad}(M_j),\quad
 r_n=r_0+\sum_{j<n}s_j.
\end{equation}
Writing $d_j=r_j/t_j$ and $\theta_j=s_j/t_j$, the exact recurrence and
its iterated form are
\[
 d_{j+1}=(d_j+\theta_j)/2,\qquad
 d_n=2^{-n}d_0+\sum_{j<n}2^{j-n}\theta_j.
\]
Every $\theta_j$ is charged from its actual new divisor; none is discarded.

\begin{theorem}[Two-step sufficient criterion]\label{tt-two-step-criterion}
If a seed and its first transform both admit \eqref{tt-profile}, with ratios
$\lambda_0,\lambda_1$, then
\[
 \frac{r_2}{t_2}\le\frac34+
       \frac{K_*}4(\lambda_0+2\lambda_1).
\]
If $\lambda_0+2\lambda_1\le1/(5K_*)$, the second transform has quality
at least $5/4$. Unbounded seed heights satisfying this condition would
disprove $abc$ at the fixed value $\varepsilon=1/8$.
\end{theorem}
\begin{proof}
The elementary estimate $r_0\le3t_0$, the bounds
$s_0\le K_*\lambda_0t_0$, $s_1\le2K_*\lambda_1t_0$, and
\eqref{tt-telescope} give
$r_2\le(3+K_*\lambda_0+2K_*\lambda_1)t_0$.
Divide by $t_2=4t_0$. Under the stated condition this is
$r_2/t_2\le4/5$. Therefore
\[
 E_{2,1/8}=4t_0-\frac98r_2
 \ge4t_0-\frac98\frac{16}5t_0=\frac25t_0.
\]
An unbounded sequence of seed heights has a subsequence on which this
fixed-epsilon excess tends to infinity.
\end{proof}

The arithmetic existence premise is open. A sequence with both
$\rho_0\to0$ and $\rho_1\to0$ would suffice, but is not supplied by
the theorem. The following analysis gives new local compatibility results
and precise restrictions on attempts to construct that sequence.

\subsection{The second norm and a fixed quartic unit}
Put $U=ab$ and retain $M_0=a^2+ab+b^2$. The second norm is
\begin{equation}\label{tt-second-norm}
 M_1=U^2+UM_0+M_0^2=c^4-Uc^2+U^2=F(a,b),
\end{equation}
where
\[
 F(X,Y)=X^4+3X^3Y+5X^2Y^2+3XY^3+Y^4.
\]
Since $0<U/M_0\le1/3$,
$M_0^2<M_1\le13M_0^2/9$. Moreover
$\gcd(M_1,abcM_0)=1$, and
\begin{equation}\label{tt-second-mod-three}
 M_1\equiv1\pmod3.
\end{equation}
For the last assertion, reduce $F$ to $(a^2+b^2)^2$ modulo three; primitivity
excludes the simultaneous zero residue and every remaining squared sum has
square one.

Let $K_0=\mathbb Q(\zeta)$ and let $\alpha$ satisfy
\[
 P(X)=X^2+(2-\zeta)X+1,
 \qquad L=K_0(\alpha).
\]
The discriminant $-1-3\zeta$ has norm thirteen, so it is not a square in
$K_0$ and $[L:\mathbb Q]=4$. The equality
$\zeta=(\alpha+1)^2/\alpha$ shows $L=\mathbb Q(\alpha)$.
Multiplying $P$ by its coefficient conjugate gives the irreducible polynomial
$F(X,1)$, with polynomial discriminant $117$: the two discriminants have
product thirteen and the squared resultant is nine. No maximal-order
identification is needed for the following argument.

\begin{lemma}[Actual unit coefficients]\label{tt-quartic-unit}
One has
\[
 M_1=N_{L/\mathbb Q}(a-\alpha b),\qquad
 (\alpha+\zeta)(\alpha+\bar\zeta)=-\bar\zeta\alpha.
\]
The element
$\kappa=(\alpha+\zeta)/(\alpha+\bar\zeta)$ is an algebraic unit outside
$K_0$, is not a root of unity, and satisfies, for $\nu=\kappa/\zeta$,
\[
 \nu+\nu^{-1}=1-3\zeta,\qquad
 \nu^4+\nu^3+9\nu^2+\nu+1=0.
\]
Finally
\begin{equation}\label{tt-linear-unit-identity}
 a-\alpha b=
 \frac{(\alpha+\zeta)\bar z-(\alpha+\bar\zeta)z}
 {\zeta-\bar\zeta}.
\end{equation}
\end{lemma}
\begin{proof}
The relative norm is
$a^2+(2-\zeta)ab+b^2=\zeta^{-1}(ab+\zeta M_0)$, whose norm to
$\mathbb Q$ is $M_1$. The quadratic relation gives the product identity.
As $\alpha$ is an algebraic unit, that identity makes both coefficient factors
units. If $\kappa\in K_0$, it cannot equal one and
$\alpha=(\zeta-\kappa\bar\zeta)/(\kappa-1)$ would belong to $K_0$,
a contradiction. Directly taking the relative norm and trace of $\kappa$
gives $\zeta^2$ and $3-2\zeta$. Division by $\zeta$ and multiplication
by the coefficient conjugate give the two assertions about $\nu$.
The nonreal value $1-3\zeta$ prevents $\nu$ from having complex absolute
value one, so neither $\nu$ nor $\kappa$ is a root of unity. Solving the
two linear equations $z=a+b\zeta$, $\bar z=a+b\bar\zeta$ proves
\eqref{tt-linear-unit-identity}.
\end{proof}

\begin{theorem}[Second-norm small-prime mass]\label{tt-second-small-primes}
Suppose the first profile is the actual pure-power representation
\[
 z=u(1+\zeta)^{e_0}w^g,\qquad e_0\in\{0,1\},\quad Q=N(w)\ge7.
\]
There is an effective absolute $D>0$, independent of $w$ and its support,
such that
\[
 v_p(M_1)\le Dp^4\log Q\log(3+g)\quad(p\mid M_1),
\]
and, for $Y\ge2$,
\[
 \log\prod_{\substack{p\mid M_1\\p\le Y}}p^{v_p(M_1)}
 \le D\log Q\log(3+g)Y^5\log Y.
\]
For $g\ge1024$, take $Y=g^{1/10}$. The ratio of this mass to $\log M_1$
is at most
$\frac D{20}\log(3+g)\log g/\sqrt g$, which tends to zero uniformly
in the size and number of norm primes of $w$.
\end{theorem}
\begin{proof}
Put $\eta=w/\bar w$ and
$\kappa_{u,e_0}=\kappa u^{-2}\zeta^{-e_0}$. Since
$z/\bar z=u^2\zeta^{e_0}\eta^g$, equation
\eqref{tt-linear-unit-identity} gives, at every place above $p\mid M_1$,
\[
 v_P(a-\alpha b)=v_P(\eta^g\kappa_{u,e_0}^{-1}-1).
\]
Every omitted factor is a unit there: $p\nmid3Q$ by
\eqref{tt-second-mod-three} and the norm-boundary gcd identity, the two
coefficient factors are global units, and $\zeta-\bar\zeta$ is supported
only above three. This includes the ramified prime thirteen.

Apply the first inequality of Theorem 1.3 of
\href{https://arxiv.org/html/2209.00275v1}{Bugeaud, \emph{B prime}},
with coefficients $g,-1$ in a field of degree at most four. The product is
not one since $\eta^g\in K_0$ and $\kappa_{u,e_0}\notin K_0$.
One has $h^*(\eta)\le\log Q$, and the finite root-of-unity twists of
$\kappa$ have fixed height. With normalized valuations $v_P(p)=1$, the
source theorem gives the stated bound for each place, up to an absolute
constant. For the rational norm one must sum the places:
\[
 v_p(M_1)=\sum_{P\mid p}f_P\operatorname{ord}_P(a-\alpha b)
 =\sum_{P\mid p}e_Pf_Pv_P(a-\alpha b),\qquad
 \sum_{P\mid p}e_Pf_P=4.
\]
Enlarge one absolute constant to obtain the assertion, including ramification.
Summation uses $\sum_{p\le Y}p^4\log p\le Y^5\log Y$.
Finally $\log M_1\ge2g\log Q$ and substitution of $Y=g^{1/10}$ give
the displayed ratio.
\end{proof}

This positive distribution theorem does not bound the second radical from
above. It forces any successful second compression to account for depth at
moving large primes.

\subsection{What fixed first roots force on second compression}
\begin{theorem}[A fixed-root bound on second pure content]\label{tt-fixed-root-content}
Fix $w$ in the preceding theorem. Let $h$ be the gcd of the positive prime
exponents of $M_1$, and choose any $p\mid Q$ with $e=v_p(Q)>0$.
If $C>2$ is the absolute constant of the established one-block two-place
estimate, then
\[
 \frac{h}{\log(3+h)}\le
 \frac{2C^2p^2}{e}(\log Q+\log4).
\]
In particular $h$ has an effective upper bound independent of $g$ and of the
unit rotation choosing positive coordinates.
\end{theorem}
\begin{proof}
By \eqref{tt-second-mod-three}, the oriented factorization of $z_1$ has no
ramified factor. Dividing all its exponents by $h$ gives an actual
representation $z_1=u_1W^h$. The boundary of the first transformed triple
contains $M_0$, hence has $p$-valuation at least $ge$.
The previously established one-block estimate applied to this triple yields,
with $t_1=\log(c^2)$,
\[
 ge\le p^2 C^2\log(3+h)\log N(W)
 \le2C^2p^2t_1\frac{\log(3+h)}h.
\]
Since $M_0\ge3c^2/4$,
$t_1\le g\log Q+\log4$. Divide by $g$ and use $g\ge1$.
Effectivity follows because $h/\log(3+h)\to\infty$.
\end{proof}

The following earlier $\rho$ estimate alone leaves a remainder window.
More generally, if
$z=u(1+\zeta)^{e_0}v_0w^{g_0}$ with fixed $w$, the established shared bound
$v_p(T_1)\le A p^2\rho_1t_1$ for a representation of the first transform
implies
\[
 \rho_1\ge
 \frac{e}{A p^2(\log Q+\lambda_0+(\log4)/g_0)}.
\]
This follows from $v_p(T_1)\ge g_0e$ and
$t_1\le g_0\log Q+\log N(v_0)+\log4$.
If $\rho_0\to0$, then $\lambda_0\to0$, $g_0\to\infty$, and
$\liminf\rho_1\ge e/(A p^2\log Q)>0$.
Taken alone, this estimate does not exclude
$\lambda_1=\rho_1/\log(4+g_1)\to0$ with a growing remainder.
The subsequent refinement in Corollary~\ref{lr-fixed-root} proves a positive
$\lambda_1$ floor when $w$ is fixed and $\lambda_0$ is bounded, closing
that asymptotic window. Neither bound has been compared with
Theorem~\ref{tt-two-step-criterion}'s finite lambda threshold, and roots
with unbounded norm remain subject to the separate necessary conditions in
Corollary~\ref{lr-moving-support}.

\subsection{An exact local obstruction and positive arbitrary-depth lifting}
\begin{theorem}[The thirteen-adic boundary]\label{tt-thirteen}
For every primitive integer pair $(a,b)$,
$13\mid F(a,b)$ if and only if $a\equiv b\pmod{13}$, and in that case
$v_{13}(F(a,b))=1$.
\end{theorem}
\begin{proof}
Modulo thirteen,
$F(a,b)=(a-b)^2(a^2+5ab+b^2)$. The second quadratic has discriminant
eight, outside the square residues $0,1,3,4,9,10,12$.
If $b=0$ modulo thirteen, primitivity gives $F=a^4\ne0$; otherwise divide
by $b^4$. This proves the equivalence. For the exact depth, write $a=b+d$:
\[
 F(b+d,b)=d^4+7bd^3+20b^2d^2+26b^3d+13b^4.
\]
At $d=13k$, this is $13b^4\not\equiv0\pmod{169}$.
\end{proof}

For the fixed root $w=2+\zeta$, exact arithmetic gives
$w^{12}=1\pmod{13}$ and $w^{10}=6+6\zeta\pmod{13}$.
Below, $F(w^g)$ denotes $F$ applied to the two integer coordinates of $w^g$.
Infinitely many exponents $g=10+12n$ have positive raw coordinates:
if $\theta=\arg w$, then $\theta/\pi$ is irrational, since a root-of-unity
quotient $w/\bar w$ would equate its distinct oriented prime ideals above
seven. The irrational rotation by $12\theta$ visits the open sector
$(0,\pi/3)$ infinitely often. Density follows from the arbitrarily small
angles in the closure of an infinite cyclic subgroup of the circle; positive
orbits and their tails have that same closure.
These actual positive pairs are primitive, because a common rational prime
divisor would force a conjugate factor into $w^g$. Their first norm is $7^g$,
while their second content is exactly one by Theorem~\ref{tt-thirteen}.
This refutes automatic pure-content transfer, not a transfer with remainders.

\begin{theorem}[Arbitrarily deep compatible contact at sixty-seven]
\label{tt-sixty-seven}
For each $k\ge1$ there is a unique class
$g=g_k\pmod{66\cdot67^{k-1}}$ within the initial class $g=1\pmod{66}$
for which $67^k\mid F(w^g)$. The classes lift compatibly, with
$g_1=1$, $g_2=199$, $g_3=9043$, and $g_4=7415893$.
For each fixed depth, infinitely many positive primitive pairs of first norm
$7^g$ have this second-norm divisibility.
\end{theorem}
\begin{proof}
Exact calculations give
\[
 F(2,1)=67,\qquad
 w^{66}\equiv3954+1541\zeta
 =1+67(59+23\zeta)\pmod{67^2}.
\]
Put $\Delta=(w^{66}-1)/67$. Modulo 67, its coordinates are $(59,23)$
and those of $w\Delta$ are $(28,61)$. Since
$\partial_aF(2,1)=91$, $\partial_bF(2,1)=86$, the exponent derivative is
$91\cdot28+86\cdot61=22\pmod{67}$.
The binomial theorem gives
\[
 (w^{66})^{67^{k-1}}\equiv1+67^k\Delta\pmod{67^{k+1}}.
\]
Inductively, all terms after the first order vanish at the required modulus
on raising to the 67th power. A Taylor expansion of the quartic therefore gives
\[
 F(w^{g+66\cdot67^{k-1}j})
 \equiv F(w^g)+22j67^k\pmod{67^{k+1}},
\]
since $w^g=w\pmod{67}$ in the initial class. Exactly one $j\pmod{67}$
kills the next digit. The first is $j=3$, yielding $g_2=199$.
Induction proves existence and uniqueness of the compatible classes.
The irrational-rotation argument above gives positive raw coordinates
infinitely often in each fixed class, and oriented factorization gives
primitivity and first norm $7^g$.
\end{proof}

At $g=199$, multiplication of the negative raw coordinates by minus one gives
a positive pair without changing the quartic. Exact division verifies that
its 337-digit second norm has valuation exactly two at 67. Its full
factorization is not asserted, and this local witness is not described as a
compressed second norm.

\begin{theorem}[Two distinct prime targets are simultaneously compatible]
\label{tt-two-prime-lifting}
For any $k,\ell\ge1$, infinitely many actual positive primitive pairs of
first norm $7^g$ satisfy $67^k967^\ell\mid M_1$.
\end{theorem}
\begin{proof}
Use the following exact local data; the derivative is along multiplication
by $(w^d-1)/p$ in the exponent direction:
\[
\begin{array}{c|c|c|c|c|c}
p&d&g_1&(w^d-1)/p\pmod p&w^{g_1}\pmod p&D_p\\\hline
67&66&21&(59,23)&(33,66)&7\\
967&966&651&(759,279)&(53,172)&39
\end{array}
\]
Both quartic residues vanish, and both derivatives are units. The preceding
binomial and Taylor argument gives a class modulo $66\cdot67^{k-1}$ and
a class modulo $966\cdot967^{\ell-1}$. Their representatives remain three
modulo six. Since $66=2\cdot3\cdot11$ and $966=2\cdot3\cdot7\cdot23$,
the two moduli have gcd exactly six for every depth pair. The generalized
Chinese remainder theorem supplies a common class, beginning with
$g=7413\pmod{10626}$ at depth one. Positive raw pairs occur infinitely
often in each common class by irrational rotation, and they remain primitive.
\end{proof}

These simultaneous local towers do not satisfy the whole-norm premise merely
by construction. Indeed Theorem~\ref{tt-second-small-primes} shows that any
fixed finite set of primes carries only $o(\log M_1)$ mass as $g\to\infty$.
The unresolved requirement is high depth across enough of the full moving
support, with the exact remainder norm below the two-step threshold.

\subsection{Exact computation and remaining scope}
The standalone integer replay \texttt{two\_step\_compatibility.py} verifies
twelve complete second-norm factorizations for the first twelve positive
rotations of $(2+\zeta)^g$. Primality is checked by trial division and each
product is reconstructed exactly. All twelve are squarefree, which proves
only that those twelve inputs do not meet the second compression threshold.
The $g=199$ witness already refutes extrapolating universal squarefreeness.
The same replay checks all relevant thirteen-adic residue cases, eight
sixty-seven-adic lifting levels, and thirty-six simultaneous depth pairs
between one and six at 67 and 967. Local checks never stand in for a missing
complete factorization.

The remaining branches are both active: fixed roots with growing second
remainders, and moving roots of unbounded norm. The new quartic unit gives
a concrete shifted residue target; the local towers prove nontrivial
compatibility; and the small-prime theorem and fixed-root content bound give
precise limitations. None establishes the unbounded family required by
Theorem~\ref{tt-half-quality} or Theorem~\ref{tt-two-step-criterion}.
